\documentclass[12pt,a4paper]{article}
\usepackage[T1]{fontenc}
\usepackage[utf8]{inputenc}
\usepackage{lmodern}
\usepackage[english]{babel}
\usepackage[margin=2.5cm]{geometry}
\usepackage{xcolor}
\usepackage{titlesec}
\usepackage{enumitem}
\usepackage{caption}
\usepackage{microtype}
% Shared by report.tex and presentation.tex.
\usepackage{graphicx}
\usepackage{booktabs}
\usepackage{tikz}
\usetikzlibrary{arrows.meta,positioning,fit,calc,shapes.misc}

\definecolor{todored}{HTML}{C0392B}
\definecolor{good}{HTML}{1E8449}
\definecolor{bad}{HTML}{B03A2E}
\definecolor{mesh}{HTML}{2E86C1}

% Anything the team still has to fill in. Search the .tex files for \todo.
\newcommand{\todo}[1]{{\color{todored}\textbf{[TODO:} #1\textbf{]}}}

% \placeholder{width}{height}{text}: red box where a figure will go.
\newcommand{\placeholder}[3]{%
  \fcolorbox{todored}{todored!5}{\parbox[c][#2][c]{\dimexpr#1-2\fboxsep-2\fboxrule\relax}%
    {\centering\color{todored}\small #3}}}

% \evidence{file}{width}{height}{what it should show}
% Shows evidence/<file> once it exists, a placeholder until then.
\newcommand{\evidence}[4]{%
  \IfFileExists{evidence/#1}%
    {\includegraphics[width=#2,height=#3,keepaspectratio]{evidence/#1}}%
    {\placeholder{#2}{#3}{missing \texttt{evidence/#1}\\[2pt]#4}}}

\tikzset{
  >={Stealth},
  svc/.style={draw, rounded corners, minimum width=2.3cm, minimum height=.75cm,
              font=\small\ttfamily, fill=white},
  hacked/.style={svc, draw=bad, fill=bad!12},
  proxy/.style={svc, fill=mesh!15, draw=mesh},
}

\usepackage[hidelinks]{hyperref}

\titlelabel{\thetitle.\quad}
\setlength{\parskip}{0.5em}
\setlength{\parindent}{0pt}

\begin{document}

% ---------------------------------------------------------------- title page
\begin{titlepage}
  \centering
  {\large Lodz University of Technology\\
  Faculty of Technical Physics, Information Technology and Applied Mathematics\\
  Institute of Information Technology\par}
  \vspace{1.5cm}
  {\large Selected Issues of Network Applications\\2026/2027\par}
  \vspace{3cm}
  {\LARGE\bfseries Modern Service Mesh: a better way to secure communication
  between microservices\par}
  \vspace{1cm}
  {\Large Stage 2 report: Proof of Concept\par}
  \vspace{3cm}
  {\large
  Jędrzej Bartoszewski, 251482\\
  Michał Domagała, 251505\\
  Kacper Maziarz, 251586\\
  Mikołaj Pawłoś, 258681\par}
  \vfill
\end{titlepage}

The problem definition, the choice of tools and the functional requirements
are unchanged from the stage 1 report. This report covers what we built
during the Proof of Concept stage and what we learned while doing it.

% ---------------------------------------------------------------- 1
\section{Work progress to date}

\begin{table}[h]
  \centering
  \begin{tabular}{clll}
    \toprule
    \# & Step & Planned & Status \\
    \midrule
    1 & k3d cluster and Online Boutique      & Weeks 3--4                 & Done    \\
    2 & Linkerd with mTLS (Proof of Concept) & Weeks 5--6, until 6.11.2026 & Done    \\
    3 & Zero-Trust authorization policies    & Week 7                     & Next    \\
    4 & Prometheus and Grafana               & Week 8                     & Planned \\
    5 & Attack simulation                    & Week 9, until 27.11.2026    & Planned \\
    \bottomrule
  \end{tabular}
  \caption{Schedule from the stage 1 report and the current status.}
\end{table}

We met the 6.11 deadline for both Proof of Concept steps, and the plan for
steps 3--5 stays as it was.

\subsection*{Step 1: cluster and test application}

The whole environment comes up with one script: \texttt{setup.sh} on Linux
and macOS, \texttt{setup.ps1} on Windows. Docker is the only thing a team
member has to install. The script downloads pinned versions of k3d (v5.9.0),
kubectl (v1.37.2) and the Linkerd CLI (edge-26.11.1) into a local
\texttt{.tools/} directory, so it needs no admin rights and ignores whatever
versions happen to be installed system-wide. It then creates a k3d cluster
named \texttt{sina} with one server and two agent nodes
(\texttt{cluster/k3d-config.yaml}), deploys Online Boutique v0.10.7 from the
manifest kept in the repo (\texttt{app/kubernetes-manifests.yaml}) and waits up
to ten minutes for every deployment to become available. Host port 8080 is
mapped to the k3d load balancer, which reaches the shop's
\texttt{frontend-external} service, so the shop is at
\url{http://localhost:8080} with no port-forward running.

Re-running the script reuses the existing cluster, which means anyone can run
it again after a failed attempt without cleaning up first.
\texttt{cluster/README.md} lists the failures we
expect people to hit: pods stuck in \texttt{Pending} or \texttt{OOMKilled}
when Docker has too little memory, Docker Desktop without WSL2 integration on
Windows, and Podman in place of Docker on Linux.

Before we added the mesh, a fresh run had the cluster ready after 27 seconds,
and all 12 pods reached \texttt{Running} less than a minute later, each with a
single container (\texttt{1/1}). Figure~\ref{fig:shop} shows the running shop.

\begin{figure}[h]
  \centering
  \evidence{00-shop.png}{\linewidth}{7cm}{Online Boutique at \texttt{localhost:8080}}
  \caption{Online Boutique running on the local k3d cluster after
    \texttt{setup.sh}, reached through the k3d load balancer on
    \texttt{localhost:8080}.}
  \label{fig:shop}
\end{figure}

\subsection*{Step 2: Linkerd and mutual TLS}

We installed Linkerd edge-26.11.1 with its CLI, and the steps live in
\texttt{setup.sh} and \texttt{setup.ps1}, so the mesh comes up with the rest of
the environment. The scripts install the Gateway API CRDs that Linkerd needs,
run \texttt{linkerd install --crds} and \texttt{linkerd install}, confirm the
result with \texttt{linkerd check}, and deploy the shop through
\texttt{linkerd inject --native-sidecar}. That command adds the
\texttt{linkerd.io/inject: enabled} annotation to each pod template as the
manifest is applied. Every Online Boutique pod then starts with a second
container, the \texttt{linkerd-proxy} sidecar, which takes over the pod's
inbound and outbound traffic. The application code stayed untouched, and the
manifest in the repo is the unmodified release file.

The \texttt{linkerd-viz} extension used for the checks below is not part of
the setup scripts yet; we installed it separately with
\texttt{linkerd viz install}. \texttt{linkerd check}, \texttt{linkerd check
--proxy} and \texttt{linkerd viz check} all passed, and all 12 pods showed
\texttt{2/2} containers.

We checked the result three ways.

\begin{enumerate}[itemsep=2pt]
  \item \texttt{linkerd viz edges} lists every connection between deployments
    and whether it's secured (Figure~\ref{fig:edges}). All 16 connections
    between application services showed as secured, and so did the 12
    connections from the \texttt{linkerd-viz} Prometheus to those services.
  \item \texttt{linkerd viz tap} streams individual live requests. Requests
    from \texttt{checkoutservice} to \texttt{paymentservice} carry
    \texttt{tls=true} (Figure~\ref{fig:tap}): all 20 captured lines did.
  \item A packet capture of \texttt{paymentservice} traffic, taken once with
    the service out of the mesh and once inside it
    (Figure~\ref{fig:capture}). We placed five test orders during each
    capture. Without the proxy, the test card number and the currency code
    were readable in the packets of every \texttt{Charge} call. With the proxy,
    neither string appeared anywhere in the capture, and the traffic ran
    between the proxies on port 4143 as unreadable data.
\end{enumerate}

The third check is the one we'll reuse in the final demo, because it shows
the difference to anyone, Kubernetes background or not.

\begin{figure}[p]
  \centering
  \evidence{04-edges.png}{\linewidth}{6cm}{\texttt{linkerd viz edges deploy -n default}}
  \caption{Connections between Online Boutique deployments as reported by
    \texttt{linkerd viz edges}. The rows for the \texttt{linkerd-viz}
    Prometheus are left out; they are secured as well.}
  \label{fig:edges}

  \vspace{1em}
  \evidence{05-tap.png}{\linewidth}{6cm}{\texttt{linkerd viz tap} from checkout to payment}
  \caption{Live requests from \texttt{checkoutservice} to
    \texttt{paymentservice}, each marked \texttt{tls=true}.}
  \label{fig:tap}
\end{figure}

\begin{figure}[p]
  \centering
  \evidence{06-before.png}{\linewidth}{8cm}{tcpdump, \texttt{paymentservice} outside the mesh}

  \vspace{1em}
  \evidence{07-after.png}{\linewidth}{8cm}{tcpdump, \texttt{paymentservice} inside the mesh}
  \caption{One packet from each capture of \texttt{paymentservice} traffic,
    taken from the node, without (top) and with (bottom) the Linkerd proxy.
    The top packet carries the test card number in the clear.}
  \label{fig:capture}
\end{figure}

% ---------------------------------------------------------------- 2
\section{Insights and challenges}

\subsection*{Encryption isn't access control}

This is the most important thing we learned in the PoC stage, and it changes
how we present the project. After step 2, every connection between meshed
services runs over mutual TLS, and each proxy knows the identity of its
caller, derived from the caller's Kubernetes ServiceAccount (Online Boutique
gives each of its 11 services its own). That stops eavesdropping and
impersonation.

It doesn't stop lateral movement. Linkerd's default inbound policy is
\texttt{all-unauthenticated}, so a meshed service still accepts requests from
any other meshed service and even from a pod outside the mesh speaking
plaintext. An attacker who controls \texttt{adservice} can't read the traffic
between checkout and payments anymore, but can still call
\texttt{paymentservice} directly and get an answer. We confirmed this with two
test pods, one outside the mesh and one inside it running under the
\texttt{adservice} ServiceAccount. Both reached the application: a
\texttt{Charge} call with an empty body came back with the payment service's
own error about the missing card, and \texttt{GetCart} on
\texttt{cartservice} succeeded. Closing that gap
is exactly what step 3 is for, and the schedule already has it in the right
order; we just hadn't appreciated how big that gap would look on a slide
until we had mTLS running.

\subsection*{Linkerd and Traefik's Gateway API CRDs}

k3s installs Traefik by default, and its CRD chart clashes with the Gateway API
CRDs that Linkerd needs: the Traefik install job kept crashing on a CRD
ownership conflict. Nothing in the project uses Traefik, because the shop is
reached through the k3d load balancer, so we disabled it in
\texttt{cluster/k3d-config.yaml}.

\subsection*{Proving encryption took more effort than enabling it}

Turning on mTLS took a few commands. Showing it to someone else took longer.
Online Boutique images are distroless, with no shell and no tcpdump, and the
pods set \texttt{runAsNonRoot: true}, which blocks the usual trick of
attaching a root debug container. We captured from the node instead, using
\texttt{kubectl debug node/\ldots{} --profile=sysadmin}.

That route has a second advantage. Inside a meshed pod, the hop from the
application to its own proxy goes over loopback and stays plaintext; a
capture there would show readable traffic and look as if mTLS had failed.
From the node you only see what leaves the pod, which is the traffic an
attacker could reach.

\subsection*{Linkerd's release model}

Since Linkerd 2.15 (February 2024), the open-source project publishes only
edge releases, and stable builds come from Buoyant, the company behind
Linkerd. We used edge release edge-26.11.1. For a student project that's fine,
but edge releases can change behaviour between versions, so our setup scripts
pin the version the same way they pin k3d and kubectl.

\subsection*{Resource cost}

Each of the 12 application pods now runs an extra proxy container, plus the
Linkerd control plane and the viz extension. In one \texttt{kubectl top}
snapshot, shortly after the pods started, the 12 proxies used 43~MiB together,
about 3.6~MiB each, next to 333~MiB for the application containers. The control
plane used 120~MiB and the viz extension 264~MiB.

The main problem on a laptop was disk space rather than memory. When the system
drive filled up on one of our Windows machines, Docker Desktop stopped
responding while the cluster kept running, so \texttt{docker info} hung and the
setup script hung with it. Restarting every deployment at once also strained
the machine: \texttt{emailservice} and \texttt{recommendationservice} failed
their health probes during start-up and restarted a few times before settling.

% ---------------------------------------------------------------- 3
\section{Impact on the next steps}

\textbf{Step 3 (Michał).} We'll switch the namespace to a default-deny
inbound policy and then allow only the calls Online Boutique actually makes,
written as Linkerd \texttt{Server}, \texttt{AuthorizationPolicy} and
\texttt{MeshTLSAuthentication} resources against the per-service identities.
\texttt{paymentservice}, for example, will accept only the
\texttt{checkoutservice} identity. Two things have to keep working: the
frontend receives traffic from outside the mesh, and the loadgenerator must
still reach it.

\textbf{Step 4 (Kacper).} \texttt{linkerd-viz} comes with its own Prometheus,
but that instance keeps data for only six hours, and Linkerd no longer bundles
Grafana. We'll start with the bundled Prometheus and add Grafana on top,
because it's already running and needs less on our laptops. If it limits us,
we'll move to a separate Prometheus that scrapes the proxies. Either way, the
goal is one Grafana dashboard where the step 5 attack shows up as denied
requests.

\textbf{Step 5 (Michał).} We'll run the same attack twice, once from a pod
outside the mesh and once from inside a compromised \texttt{adservice}, both
calling \texttt{paymentservice}. Before step 3 both calls succeed; afterwards
both should fail with gRPC \texttt{PermissionDenied}.

\end{document}
